\documentclass[10pt,a4paper]{amsart}
\usepackage{amsmath,amssymb,amsthm,graphicx}
\usepackage{hyperref}
\newtheorem{thm}{Theorem}[section]
\newtheorem{thm*}{Theorem}
\newtheorem{lem}[thm]{Lemma}
\newtheorem{prop}[thm]{Proposition}
\newtheorem{cor}[thm]{Corollary}
\newtheorem{claim}[thm]{Claim}
\newtheorem{defprop}[thm]{Definition-Proposition}
\newtheorem{conj}[thm]{Conjecture}
\newtheorem{fact}[thm]{Fact}

\newtheorem{introthm}{Theorem}
\newtheorem{introprop}{Proposition}

\renewcommand{\theintrothm}{\Alph{introthm}}
\renewcommand{\theintroprop}{\Roman{introprop}}

{\theoremstyle{definition}

\newtheorem{df}[thm]{Definition}
\newtheorem{ex}[thm]{Example}
\newtheorem{qu}[thm]{Question}
\newtheorem{rmk}[thm]{Remark}
\newtheorem{nota}[thm]{Notation}
\newtheorem{assump}{Assumption}}

\def\nn{\nonumber}


\def\calB{{\mathcal B}}
\def\calP{{\mathcal P}}
\def\M{{\mathcal M}}
\def\Mco{{{}^{\rm c/o}\mathcal M}_{g,\delta_1,\dots,\delta_n}^{s,\beta}}
\def\Mcos{{{}^{\rm c/o}\mathcal M}_{g,\delta_1,\dots,\delta_n}^{s_1\dots s_k,\beta}}
\def\Sull{\rm {\mathcal S}ull^{\rm c/o}}
\def\wSull{\widetilde{\rm {\mathcal S}ull}^{\rm c/o}}
\def\Rp{{\mathbb R}_{>0}}
\def\la{\langle}
\def\ra{\rangle}

%\def\empty{\varnothing}

\def\idagger{r^{\dagger}}
\def\rdagger{r^{\dagger}}
\def\res{r}
\def\oto{\otimes \dots \otimes}
\newcommand{\aoa}[2]{a_{#1}\otimes \dots \otimes a_{#2}}
\newcommand{\bob}[2]{b_{#1}\otimes \dots \otimes b_{#2}}
\newcommand{\ata}[2]{a_{#1}\cdots a_{#2}}
\newcommand{\aca}[2]{a_{#1}, a_{#2}}
\newcommand{\set}[2]{\{#1,\dots, #2\}}

\def\wt{\omega}

\def\CO{{\mathcal O_\mu}}
\def\calO{{\mathcal O}}
\def\a{\alpha}
\def\t{\tau}
\def\B{\mathcal B}
\def\Z{\mathbb Z}
\def\R{\mathbb R}
\def\ba{\mathbf{a}}
\def\d{d^{{\rm CH}^*}}
\def\s{s^{{\rm CH}^*}}

%% Next two commands enlarge space for figures on the top of a page
\renewcommand{\textfraction}{0.01}
\renewcommand{\topfraction}{0.99}
\title{A CW model for the open/closed Sullivan spaces: cells indexed by Sullivan-type arc graphs, cellular gluing maps and a deformation retraction onto the subspace of unit incoming weights}
\author{Ralph M. Kaufmann}
\date{doi:10.3842/SIGMA.2010.036; CC BY 4.0}
\begin{document}
\maketitle
\setcounter{section}{6}
\setcounter{thm}{2}
\noindent\textit{Editorial note.} All delivered mathematical text is the author's own, copied from the fixed publisher source (SIGMA 6 (2010) 036, statement source lines 1575-1580 and proof source lines 1582-1611); no statement, proof step or formula was rewritten, shortened or removed.\par\medskip
\noindent\textit{Editorial note.} The publisher class sigma is not available here and is replaced by a plain amsart wrapper; the preamble is the author's own theorem declarations and macro block copied verbatim from source lines 24-90, with no mathematical macro invented, renamed or added.\par\medskip
\noindent\textit{Editorial note.} The excerpt carries no section heading, so the section counter is fixed at 6 and the author's shared theorem counter at 2 before the body; the counted statement therefore prints as Proposition 6.3, which is its number in the article under the author's own declaration newtheorem\{thm\}\{Theorem\}[section]. The only other material before it is the context paragraph, which defines the object and carries no theorem number.\par\medskip
\noindent\textit{Editorial note.} The proof cites the author's earlier cell-operad paper once, at source line 1607, for the associativity of the cellular gluing maps in the closed case. The complete bibliography of the article is therefore carried verbatim as a separate source file (source lines 2295-2454, copied without change) and hyperref is loaded so that its href entries compile; no bibliography entry was added, dropped or edited.\par\medskip
\noindent\textit{Editorial note.} The paragraph at source lines 1568-1573 is carried as context because it introduces Sull\_1, the object the statement is about; it is not counted as part of the proof.\par\medskip
\noindent\textit{Editorial note.} No figure or table is referenced by the retained spans, so none is carried.\par\medskip
\noindent\textit{Allowed standard prerequisites (editorial).}\begin{itemize}
\item Weighted arc families, their alpha weights and the Sullivan-type condition. An element of the open/closed Sullivan space wSull(n\_1,n\_2,m\_1,m\_2) is a projectivized (PMC) orbit of an arc family on a brane-labelled windowed surface whose arcs run from incoming to outgoing windows; the alpha weight of a window is the sum of the weights of the arcs incident to it.. the article's own review of the model at source lines 250-475, in particular the arc-family and Sullivan conventions. Only the notion of the weight of a window, of an arc graph of Sullivan type and of the PMC action is assumed; the cell decomposition, the retraction and the cellularity of the gluings are the retained obligations.
\item Standard cell-complex language: cells, interiors of cells, deformation retraction, cellular maps. A CW complex is a Hausdorff space with a filtration by closed cells whose interiors are disjoint open cells and whose attaching maps land in lower skeleta; a deformation retraction is a homotopy from the identity to a retraction fixing the subspace.. standard topology; used at source lines 1583-1605. The vocabulary in which the proposition is stated. The author's argument that the subspace really carries such a structure, with the stated cells and attaching maps, is retained in full.
\item Cell operad model for the closed sector of the same author's earlier work. The cellular chains of the closed analogue form a PROP cell model, and the associativity of the cellular gluing maps on the nose is established there.. cited by the author as [hoch1, Theorem 5.33] at source line 1607; the entry is carried in the bibliography supplied with this excerpt. This is a pointer to the closed half of the argument that the author transfers; the open half, where punctures are kept and the boundary twist condition is absent, is the retained obligation (source lines 1607-1610).
\end{itemize}\medskip
This paragraph is an application of the methods set forth in \cite{hoch1}.
We def\/ine the following subspaces of $\wSull(n_1,n_2,m_1,m_2)$  we let
${\Sull}_1(n_1,n_2,m_1,m_2)$ be the subspace of all $\alpha\in\wSull(n_1,n_2,m_1,m_2)$
such that the {\em $\alpha$ weight of each ``in'' window is $1$}.

\begin{prop}
${\Sull}_1(n_1,n_2,m_1,m_2)$ is a CW complex, is a sub-topological quasi-PROP and is a deformation retract of $\wSull(n_1,n_2,m_1,m_2)$.
The dimension $k$ cells of this complex are indexed by arc graphs of Sullivan type with $k+1$ arcs
and their attaching maps are given by deleting arcs and identifying this boundary with the cell of lower
dimension.
\end{prop}

\begin{proof}
The fact that this is a deformation retract is shown again by using a f\/low. This time the f\/low
scales the weights of the arcs incident to an ``in'' window  $w$ by a factor of $1-t(1/\alpha(w)-1)$.

Since for a graph of Sullivan type, all the arcs can be enumerated
by going along the ``in'' windows, we see that before taking PMC
orbits, the graphs of a given Sullivan type with $\alpha$ weight on
the ``in'' windows is simply a product of simplices. Now there is no
PMC isotropy inside this product of simplices, but under the action
sides may become identif\/ied. So for each arc graph $\Gamma$  of
Sullivan type on a brane labelled windowed surface $F$, we obtain a
cell $C(\Gamma)$ whose interior $\dot C(\Gamma)$ is given by a
product of open simplices $\dot\Delta^k$:
\begin{equation*}
\dot  C(\Gamma) =\prod_{w \in \{\text{ ``in'' windows of $F$\}}} \dot\Delta^{|\{\text{arcs incident to $w$}\}|}.
\end{equation*}
The statement about the attaching maps follows directly from the topology in $\wSull$,
where we simply delete
an arc in the limit where its weight goes to zero.
The fact that this is a sub-topological quasi PROP is immediate upon
noticing that the property that the $\alpha$ weight of each ``in'' window is  one is stable under the operation of gluing.

Since the gluing a window with $n$ arcs to a window with $m$ arcs produced at most $n+m-1$ arcs
we see that the gluing maps are indeed cellular and there are induced maps on the cellular level.

It remains to prove that these maps are associative on the nose. This follows in the same way as in \cite[Theorem~5.33]{hoch1}.
The proof there essentially goes over to the current situation for the closed part. The open part is actually simpler,
since we do not have to worry about the condition of ``twisted at the boundary'', since we keep the punctures
upon gluing.
\end{proof}

\begin{thebibliography}{99}

\footnotesize\itemsep=0pt

\bibitem{ramirez}
Baas N.A., Cohen  R.L., Ram\'{\i}rez A.,
The topology of the category of open and closed strings,
in Recent Developments in Algebraic Topology,
{\it Contemp. Math.}, Vol.~407, Amer. Math. Soc., Providence, RI, 2006,  11--26,
\href{http://arxiv.org/abs/math.AT/0411080}{math.AT/0411080}.

\bibitem{BCT}
Blumberg A.J., Cohen R.L., Teleman C.,
Open-closed f\/ield theories, string topology, and Hochschild homo\-lo\-gy,
in Alpine Perspectives on Algebraic Topology, Editors C.~Ausoni, K.~Hess and J.~Scherer,
{\it Contemp. Math.}, Vol.~504, Amer. Math. Soc., Providence, RI, 2009, 53--76,
\href{http://arxiv.org/abs/0906.5198}{arXiv:0906.5198}.

%\bibitem{CJ}
%Cohen R.L., Jones J.D.S.,
%A homotopy theoretic realization of string topology,
%\href{http://dx.doi.org/10.1007/s00208-002-0362-0}{{\it Math. Ann.}} {\bf 324} (2002), 773--798,
%\href{http://arxiv.org/abs/math.GT/0107187}{math.GT/0107187}. !!!Not cited!!!

 \bibitem{CS}
 Chas M., Sullivan D.,
String topology,
{\it Ann. of Math.}, to appear,
\href{http://arxiv.org/abs/math.GT/9911159}{math.GT/9911159}.

\bibitem{godin}
Godin V.,
Higher string topology operations,
\href{http://arxiv.org/abs/0711.4859}{arXiv:0711.4859}.

\bibitem{HVZ}
Harrelson E., Voronov A.A., Zuniga J.J.,
Open-closed moduli spaces and related algebraic structures,
\href{http://arxiv.org/abs/0709.3874}{arXiv:0709.3874}.

\bibitem{J}
Jones J.D.S.,
Cyclic homology and equivariant homology,
\href{http://dx.doi.org/10.1007/BF01389424}{{\it Invent. Math.}} {\bf 87} (1987), 403--423.


%\bibitem{KLi1}
%Kapustin A., Li Y.,
%D-branes in Landau--Ginzburg models and algebraic geometry,
%\href{http://dx.doi.org/10.1088/1126-6708/2003/12/005}{{\it J. High Energy Phys.}} {\bf 2003} (2003), no.~12, 005, 44~pages,
%\href{http://arxiv.org/abs/hep-th/0210296}{hep-th/0210296}. !!!Not cited!!!

%\bibitem{KR}
%Kapustin A., Rozansky L.,
%On the relation between open and closed topological strings,
%\href{http://dx.doi.org/10.1007/s00220-004-1227-z}{{\it  Comm. Math. Phys.}} {\bf  252} (2004), 393--414,
%\href{http://arxiv.org/abs/hep-th/0405232}{hep-th/0405232}. !!!Not cited!!!

\bibitem{hoch1}
Kaufmann R.M.,
Moduli space actions on the Hochschild co-chains of a Frobenius algebra. I.~Cell operads,
\href{http://dx.doi.org/10.4171/JNCG/10}{{\it J. Noncommut. Geom.}} {\bf 1} (2007), 333--384,
\href{http://arxiv.org/abs/math.AT/0606064}{math.AT/0606064}.

\bibitem{hoch2}
Kaufmann R.M.,
Moduli space actions on the Hochschild co-chains of a Frobenius algebra. II.~Correlators,
\href{http://dx.doi.org/10.4171/JNCG/22}{{\it J. Noncommut. Geom.}} {\bf 2} (2008), 283--332,
\href{http://arxiv.org/abs/math.AT/0606065}{math.AT/0606065}.


\bibitem{romp}
Kaufmann R.M.,
Noncommutative aspects of open/closed strings via foliations,
\href{http://dx.doi.org/10.1016/S0034-4877(08)80016-X}{{\it Rep. Math. Phys.}} {\bf 61} (2008), 281--293,
\href{http://arxiv.org/abs/0804.0608}{arXiv:0804.0608}.

\bibitem{postnikov}
Kaufmann R.M.,
Dimension vs. genus: a surface realization of the little $k$-cubes and an $E_{\infty}$ operad,
in Algebraic Topology -- Old and New (M.M.~Postnikov Memorial Conference),
\href{http://dx.doi.org/10.4064/bc85-0-17}{{\it Banach Center Publ.}}, Vol.~85, Polish Acad. Sci., Warsaw, 2009, 241--274,
\href{http://arxiv.org/abs/0801.0532}{arXiv:0801.0532}.

\bibitem{KP}
Kaufmann R.M., Penner R.B.,
Closed/open string diagrammatics,
\href{http://dx.doi.org/10.1016/j.nuclphysb.2006.03.036}{{\it Nuclear Phys.~B}} {\bf 748} (2006), 335--379,
\href{http://arxiv.org/abs/math.GT/0603485}{math.GT/0603485}.

\bibitem{KLP}
Kaufmann R.M., Livernet  M., Penner R.B.,
Arc operads and arc algebras,
\href{http://dx.doi.org/10.2140/gt.2003.7.511}{{\it Geom. Topol.}} {\bf 7} (2003), 511--568,
\href{http://arxiv.org/abs/math.GT/0209132}{math.GT/0209132}.

\bibitem{LM}
Losev A., Manin Yu.,
New moduli spaces of pointed curves and pencils of f\/lat connections,
\href{http://dx.doi.org/10.1307/mmj/1030132728}{{\it Michigan Math.~J.}} {\bf 48}  (2000), 443--472,
\href{http://arxiv.org/abs/math.AG/0001003}{math.AG/0001003}.
%Dedicated to William Fulton on the occasion of his 60th birthday.

\bibitem{MSS}
 Markl M., Shnider S., Stashef\/f J.,
Operads in algebra, topology and physics,
{\it Mathematical Surveys and Monographs}, Vol.~96, Amer. Math. Soc., Providence, RI, 2002.


%\bibitem{Merk}
%Merkulov S.A.,
%De Rham model for string topology,
%\href{http://dx.doi.org/10.1155/S1073792804132662}{{\it  Int. Math. Res. Not.}} {\bf  2004} (2004),  no.~55, 2955--2981,
%\href{http://arxiv.org/abs/math.AT/0309038}{math.AT/0309038}. !!!Not Cited in the paper!!!


\bibitem{P1}
Penner R.C.,
Decorated Teichm\"uller theory of bordered surfaces,
\href{http://projecteuclid.org/getRecord?id=euclid.cag/1098468019}{{\it Comm. Anal. Geom.}} {\bf  12}  (2004),  793--820,
\href{http://arxiv.org/abs/math.GT/0210326}{math.GT/0210326}.

\bibitem{Qu}
Quillen D.,
Elementary proofs of some results of cobordism theory using Steenrod operations,
\href{http://dx.doi.org/10.1016/0001-8708(71)90041-7}{{\it Adv. Math.}} {\bf 7} (1971), 29--56.

\bibitem{St}
Strebel K.,
Quadratic dif\/ferentials,
{\it Ergebnisse der Mathematik und ihrer Grenzgebiete (3)}, Vol.~5,
Springer-Verlag, Berlin, 1984.

\bibitem{C1}
Sullivan D.,
Sigma models and string topology,
in  Graphs and Patterns in Mathematics and Theoretical Physics,
{\it Proc. Sympos. Pure Math.}, Vol.~73, Amer. Math. Soc., Providence, RI, 2005,  1--11.


\bibitem{C2}
Sullivan D.,
Open and closed string f\/ield theory interpreted in classical algebraic topology,
in  Topology, Geometry and Quantum Field Theory,  \href{http://dx.doi.org/10.1017/CBO9780511526398.014}{{\it London Math. Soc. Lecture Note Ser.}}, Vol.~308,
Cambridge Univ. Press, Cambridge, 2004,  344--357,
\href{http://arxiv.org/abs/math.QA/0302332}{math.QA/0302332}.


\bibitem{sullsurv}
Sullivan D.,
String topology: background and present state,
\href{http://arxiv.org/abs/0710.4141}{arXiv:0710.4141}.

\bibitem{TZ}
Tradler T., Zeinalian M.,
Algebraic string operations,
\href{http://dx.doi.org/10.1007/s10977-007-9005-2}{{\it $K$-Theory}} {\bf 38} (2007), 59--82,
\href{http://arxiv.org/abs/math.QA/0605770}{math.QA/0605770}.

\end{thebibliography}
\end{document}
